%% kapa: exam=1 round=33 year=2022 period=1 form=A subject=real_estate subject_ko=부동산학원론 seq=39 number=119
%% source: Q-Net 감정평가사 기출문제 (cst00309) artlSeq=5209997 file=raw/1008/5209997/2022년도 제33회 감정평가사 1차시험 1교시(민법,경제학,부동산학).pdf
%% conversion: kordoc 4.24.51 -> tools/gen1.py (auto); review=proofread-sonnet; flags=none
\begin{question}{119}
감정평가에 관한 규칙에서 규정하고 있는 내용으로 옳지 않은 것은?
\begin{choices}
  \item 기업가치의 주된 평가방법은 수익환원법이다.
  \item 적정한 실거래가는 감정평가의 기준으로 적용하기에 적정하다고 판단되는 거래가격으로서, 거래시점이 도시지역은 5년 이내, 그 밖의 지역은 3년 이내인 거래가격을 말한다.
  \item 시산가액 조정 시, 공시지가기준법과 그 밖의 비교방식에 속한 감정평가방법은 서로 다른 감정평가방식에 속한 것으로 본다.
  \item 필요한 경우 관련 전문가에 대한 자문등을 거쳐 감정평가할 수 있다.
  \item 항공기의 주된 평가방법은 원가법이며, 본래 용도의 효용가치가 없는 물건은 해체처분가액으로 감정평가할 수 있다.
\end{choices}
\end{question}
\answer{2}
